\section{Endpoint-modular raw laws and likelihoods}
\label{sec:v55-endpoint-laws}
The error-free positive-source estimate gives more than a larger open
interval of powers. It yields quantitative local norms of each
physical remainder at every finite count, and a logarithmically
weakened endpoint law. We keep the arithmetic transition kernel in
all these conclusions, including through a change of arithmetic type.

\subsection{Small mass and a weak endpoint}
\begin{lemma}[Interpolation from the distribution tail]
\label{lem:v55-weak-interpolation}
Let $f\ge0$ on a finite measure space satisfy
$\int f\le\delta$ and
$|\{f>L\}|\le C_0L^{-p_c}$ for $L>0$. For $1<q<p_c$,
\begin{equation}\label{eq:v55-weak-interpolation}
 \int f^q\le C_{q,p_c,C_0}
                 \delta^{(p_c-q)/(p_c-1)}.
\end{equation}
\end{lemma}
\begin{proof}
Markov's inequality supplies the other bound
$|\{f>L\}|\le\delta/L$. For $\delta>0$, split
$q\int_0^\infty L^{q-1}|\{f>L\}|\,dL$ at
$T=(C_0/\delta)^{1/(p_c-1)}$. The lower part is bounded by
$q\delta T^{q-1}/(q-1)$ and the upper part by
$qC_0T^{q-p_c}/(p_c-q)$. Both have the exponent in
\eqref{eq:v55-weak-interpolation}. If $\delta=0$, the assertion is
immediate. This proof applies to a normalized finite measure and does
not require it to be a probability.
\end{proof}

Write $B_{\varepsilon,j}=m^2b^{\varepsilon,j,1}_{n,k,m,R}$ for
$j=\mathrm{inc},\mathrm{clr}$, and include $j=\mathrm{phys}$ for
their sum. For a complex insertion $w$ of modulus at most $M$, use
$B_{\varepsilon,j}^w$ for the corresponding density.
\begin{theorem}[Finite-count norms of both physical remainders]
\label{thm:v55-physical-remainders}
For $1<q<p_c$, $0<\varepsilon<\varepsilon_0$ and every $m\ge2$,
\begin{equation}\label{eq:v55-physical-Lq}
 \sup_{R,n,k,t}\int_t^{t+h}|B_{\varepsilon,j}^w(u)|^q\,du
       \le C_q M^q(1+h)\varepsilon^{9(p_c-q)}.
\end{equation}
The same inequality, with a changed constant, holds for
$B_{\varepsilon,j}^w-K_B*B_{\varepsilon,j}^w$, uniformly over all
reconstruction bandwidths $B>0$. The case $q=1$ has the bound
$CM(1+h)\varepsilon^{1/16}$.
\end{theorem}
\begin{proof}
Source domination gives
$|B_{\varepsilon,j}^w|\le M B_{\varepsilon,j}\le MP$.
On the measure $du/(1+h)$ restricted to the interval,
Theorem~\ref{thm:v55-error-free-mass} gives first moment
$C\varepsilon^{1/16}$, while
Corollary~\ref{cor:v55-physical-endpoint} gives the weak endpoint.
Apply Lemma~\ref{lem:v55-weak-interpolation} and use the identity
\[
 \frac1{16}\frac{p_c-q}{p_c-1}=9(p_c-q).
\]
The case $M=0$ is immediate; otherwise divide the inserted density
by $M$. For convolution, Minkowski's integral inequality and translation
of the roof interval give
\[
 \|K_B*B_{\varepsilon,j}^w\|_{L^q(I)}
 \le\int|K_B(v)|\,
                 \|B_{\varepsilon,j}^w\|_{L^q(I-v)}\,dv.
\]
The kernel has $\|K_B\|_1=\|K_1\|_1$. This proves the all-band
assertion without a spectral statement at that band. The first-moment
case follows by the same argument. All exact labels and the unchanged
roof value remain in these densities.
\end{proof}

\subsection{A logarithmically weakened endpoint}
For $\beta>1$ and $x\ge0$ define a convex function by
\begin{equation}\label{eq:v55-young-function}
 \Phi_\beta(x)=\int_0^x
       \frac{(x-t)t^{p_c-2}}{[\log(e+t)]^\beta}\,dt.
\end{equation}
It has $\Phi_\beta(0)=\Phi_\beta'(0)=0$ and
$\Phi_\beta''(x)=x^{p_c-2}/[\log(e+x)]^\beta>0$ for $x>0$.
Its behavior is comparable to $x^{p_c}$ near zero and to
$x^{p_c}/[\log(e+x)]^\beta$ at infinity, with constants depending
on $\beta$. To check the latter assertion, integrate over
$t\in[x/4,x/2]$ for the lower bound. For the upper bound split at
$\sqrt x$: the first part is $O(x^{(p_c+1)/2})$ and the second at
most $C x^{p_c}/(\log x)^\beta$. These also imply the doubling
bound $\Phi_\beta(Cx)\le C'\Phi_\beta(x)$ for each fixed $C$.

\begin{lemma}[Uniform endpoint modular and small-mass convergence]
\label{lem:v55-modular}
Suppose a family $f\ge0$ has uniformly bounded first moments and
weak $L^{p_c}$ constants. Then its $\Phi_\beta$ integrals are
uniformly bounded for every $\beta>1$, and
\begin{equation}\label{eq:v55-modular-tail}
 \int_{\{f>T\}}\Phi_\beta(f)
             \le C_\beta[\log(e+T)]^{1-\beta}\qquad(T\ge2).
\end{equation}
If also $\int f\le\delta\le e^{-2}$, then
\begin{equation}\label{eq:v55-small-modular}
 \int\Phi_\beta(f)\le C_\beta\left(
  \delta^{1/2}+[1+\log(1/\delta)]^{1-\beta}\right).
\end{equation}
\end{lemma}
\begin{proof}
The distribution formula for the truncated integral is
\[
 \int_{\{f>T\}}\Phi_\beta(f)
 =\Phi_\beta(T)|\{f>T\}|+
       \int_T^\infty\Phi_\beta'(L)|\{f>L\}|\,dL.
\]
The same split at $\sqrt L$ used above gives
$\Phi_\beta'(L)\le C_\beta
 L^{p_c-1}/[\log(e+L)]^\beta$ for $L\ge2$.
The weak endpoint bounds the last integrand by
$C_\beta/(L(\log L)^\beta)$, whose integral proves
\eqref{eq:v55-modular-tail}. The part $f\le2$ is controlled by the
first moment. For the small-mass assertion put
$T=\delta^{-1/(2(p_c-1))}$. On $f\le T$ use
$\Phi_\beta(f)\le C f^{p_c}\le CT^{p_c-1}f$; its integral is at
most $C\delta^{1/2}$. Apply \eqref{eq:v55-modular-tail} to the rest.
\end{proof}

\begin{theorem}[Explicit subcritical and endpoint-modular raw laws]
\label{thm:v55-raw-endpoint}
For every fixed $h>0$, the complete scalar raw law holds in local
$L^q$ for every $1<q<145/144$. It also satisfies, for every $\beta>1$,
\begin{equation}\label{eq:v55-raw-modular}
 \sup_{R,n,k,t}\frac1h\int_t^{t+h}
              \Phi_\beta(|P(u)-G(u)|)\,du\longrightarrow0.
\end{equation}
Here $G(u)=\mathcal L_{m,R}(k_1,k_2,u,n)$, with no positivity
assumption. The same conclusions hold with $|P-G|$ replaced by
\begin{equation}\label{eq:v55-path-error}
 E^{\rm path}(u)=
 \|\mathbf P(u)-G(u)\mathsf W_R\|_{\mathrm{BL}^*}.
\end{equation}
For the actual-return path, the assertion is on the central target
classes of Lemma~\ref{lem:v52-return-transfer}.
\end{theorem}
\begin{proof}
The transition kernel is uniformly bounded. Hence
$|P-G|\le P+C$ has uniformly bounded weak $L^{p_c}$ constants on
translated intervals. Its local first moment tends uniformly to zero
by Theorem~\ref{thm:v49-raw-local-TV}. Apply
Lemmas~\ref{lem:v55-weak-interpolation} and \ref{lem:v55-modular}.
The path error in \eqref{eq:v55-path-error} is at most $P+|G|$,
and its first moment tends to zero by
Theorem~\ref{thm:v52-path-remainder}. The same proof applies; it
uses the common roof versions constructed there, not a null set chosen
separately for each path test. The return-path first-moment transfer
is precisely Lemma~\ref{lem:v52-return-transfer}. Its source remains
positive, so the same scalar domination applies on the stated central
class, with its prescribed enlargement across the roof window.
The constants may depend on the fixed $h$, which is never allowed to
shrink with $m$ in this argument.
\end{proof}

For example, if the corresponding uniform mean error is $e_m\to0$,
\eqref{eq:v55-small-modular} gives the explicit modulus
$C_\beta(e_m^{1/2}+(1+|\log e_m|)^{1-\beta})$ for large $m$.
This is a modulus in the known error, not an asserted collision-count
rate. The normalized positive physical remainders separately satisfy
\[
 \sup_{R,n,k,t,m}\frac1{1+h}\int_t^{t+h}
       \Phi_\beta(B_{\varepsilon,j}(u))\,du
       \le C_\beta(1+|\log\varepsilon|)^{1-\beta}
\]
for sufficiently small $\varepsilon$, by the same lemma. Arbitrary
bounded insertions obey the corresponding bound by domination and
the doubling property. This endpoint statement is still an integral
estimate and does not remove their possible positive essential heights.

\subsection{The original conditional roof likelihood}
\begin{theorem}[Forward likelihood up to the explicit endpoint]
\label{thm:v55-endpoint-likelihood}
On a fixed roof window $I$ of length $h>0$, assume the pointwise
arithmetic floor $G\ge d>0$. Put
\[
 F=\int_I P,\qquad H=\int_I G,\qquad
 d\pi=P\,du/F,\quad dQ=G\,du/H,\quad r=d\pi/dQ.
\]
Uniformly in the same exact targets, for all $1<q<145/144$ and
$\beta>1$,
\begin{equation}\label{eq:v55-likelihood-modular}
 \int|r-1|^q\,dQ\longrightarrow0,\qquad
 \int\Phi_\beta(|r-1|)\,dQ\longrightarrow0.
\end{equation}
In particular $D(\pi\Vert Q)\to0$ and
$D_q(\pi\Vert Q)\to0$ for every fixed $1<q<145/144$.
The second integral in \eqref{eq:v55-likelihood-modular} is the
convex $f$-divergence generated by
$f(x)=\Phi_\beta(|x-1|)$.
\end{theorem}
\begin{proof}
The scalar first-moment law gives $F-H\to0$. The floor and the
upper bound on $G$ imply $dh\le H\le Ch$ and, for all sufficiently
large counts, $F\ge dh/2$. The exact likelihood satisfies
\[
 r-1=\frac{H(P-G)}{FG}+\frac{H-F}{F}.
\]
The first term has the subcritical and endpoint-modular convergence
of Theorem~\ref{thm:v55-raw-endpoint}; the reference density $G/H$
is uniformly bounded on $I$. For the modular use convexity in the
form $\Phi_\beta(x+y)\le\frac12\Phi_\beta(2x)+
\frac12\Phi_\beta(2y)$ and the fixed-factor doubling bound.
The constant term tends to zero. This proves
\eqref{eq:v55-likelihood-modular}. Convexity and monotonicity of
$\Phi_\beta$ show that $f$ is convex and nonnegative with $f(1)=0$.
For relative entropy use
$0\le x\log x-x+1\le C_q|x-1|^q$ with $1<q<145/144<2$.
For R\'enyi divergence, $L^q(Q)$ convergence gives
$\int r^q\,dQ\to1$. These are the unchanged conditional laws;
a positive integrated reference mass without the pointwise floor
would not justify the likelihood calculation.
\end{proof}

\paragraph{Proof dependencies and pointwise scope.}
The new route consists of the exact geometric global strip mass,
one fixed-band finite-count local strip estimate, and the polynomial
protected height. It uses no complete exponential height cap, no new
high-frequency resolvent, and no extension through grazing or a
competing hit. Thus it is not an invocation of the billiard-flow mixing
result of \cite{BDL} as an exact occupation--roof density theorem, nor
a differentiation of the suspension interval theorem of \cite{DN}.
The density and its arithmetic main term are unchanged. The positive
height criterion of Corollary~\ref{cor:v51-positive-criterion}, and
hence the uniform pointwise same-roof bridge implication, remain
separate from all finite-norm conclusions proved here.
